\section{Question 2: Why Is Behavior Unchanged at $\mu = 0.5$?}

At $\mu = 0.5$ the weighted utility~\eqref{eq:weighted_util} becomes:
%
\begin{align}\label{eq:mu_half}
    U\big|_{\mu=0.5}
      &= \tfrac{1}{2}\,u_1 + \tfrac{1}{2}\,u_2
       = \tfrac{1}{2}(u_1 + u_2)
       = \tfrac{1}{2}\,U_{\text{baseline}}
\end{align}
%
so the weighted objective is a positive scalar multiple of the baseline utility in equation~\eqref{eq:baseline_util}.

The key observation is that utility is \emph{ordinal}: optimal choices depend only on the \emph{ranking} of alternatives, not on the absolute values of $U$. Multiplying the objective by any positive constant $\alpha$ leaves the argmax unchanged, because
\[
    \argmax_{h_1,h_2}\; \alpha\, U = \argmax_{h_1,h_2}\; U \quad \text{for all } \alpha > 0.
\]
Since $U|_{\mu=0.5} = \tfrac{1}{2}\,U_{\text{baseline}}$, the optimizer chooses identical hours $h_1$ and $h_2$ in every period, and therefore all simulated labor-market outcomes are unchanged.

This invariance is special to $\mu = 0.5$. Away from it the weighting is no longer a uniform rescaling: it changes the \emph{relative} disutility of the two members' labor, so the household tilts hours toward the member whose effort is weighted less. As $\mu$ rises above $0.5$, member~1's labor disutility receives more weight and member~1 works less while member~2 works more, and symmetrically for $\mu < 0.5$. Equal weights are thus the unique point at which the planner's problem coincides with the unweighted baseline.